% Keyboard Navigation Implementation

This document presents the technical implementation of keyboard navigation features developed to enhance accessibility within the application's editor and app builder environments. The implementation enables comprehensive keyboard-based interaction with UI components, ensuring compliance with WCAG 2.1 Level AA accessibility standards.

\subsubsection{Component Accessibility}

\textbf{Interactive Components (Button, Toggle, Link)}

For Button, Toggle, Link, and similar interactive components, we implemented a consistent keyboard handling pattern:

\begin{lstlisting}[caption=Interactive Components Keyboard Pattern, label=lst:button-toggle-link-focus]
// frontend/src/Editor/Components/Button.jsx, ToggleV2.jsx, Link.jsx, etc.
const handleKeyDown = (e) => {
  if (e.key === 'Enter' || e.key === ' ') {
    e.preventDefault();
    fireEvent('onClick'); // or handleToggleChange(), handleLinkClick()
  }
};

<component
  onKeyDown={handleKeyDown}
  tabIndex={disabledState ? -1 : 0}
  role="button" // or "switch" for toggles, "link" for links
  aria-label={label}
/>
\end{lstlisting}

This pattern employs both Enter and Space keys for activation, dynamically manages tab order based on disabled state (\texttt{tabIndex=\{0\}} includes the element in keyboard navigation, \texttt{tabIndex=\{-1\}} removes disabled elements), and uses appropriate ARIA roles (\texttt{button}, \texttt{switch}, \texttt{link}) to communicate component purpose to assistive technologies.

\textbf{Dropdown Component}

Dropdown menus leverage React-Select's built-in accessibility features while extending keyboard navigation capabilities for seamless integration with the application's navigation system.

\begin{lstlisting}[caption=Dropdown Keyboard Support, label=lst:dropdown-focus]
// frontend/src/Editor/Components/DropdownV2/DropdownV2.jsx
const handleKeyDown = (e) => {
  if (e.key === 'Enter' && !isMenuOpen) {
    e.preventDefault();
    setIsMenuOpen(true);
  }
};

<Select
  onKeyDown={handleKeyDown}
  styles={{ control: (base) => ({ tabIndex: disabledState ? -1 : 0 }) }}
  components={{ MenuList: CustomMenuList }}
  aria-label={label}
/>
\end{lstlisting}

React-Select provides inherent keyboard support for arrow key navigation, Enter to select, and Escape to close. The custom \texttt{CustomMenuList} component ensures menu items are automatically focusable through React-Select's built-in option management, with each option receiving proper ARIA attributes and keyboard event handlers. When the menu opens, focus automatically moves to the search input (if enabled) or the first menu item, enabling immediate arrow key navigation through available options.

\textbf{Input Field Component}

Text input fields implement keyboard accessibility through semantic HTML elements with proper ARIA attributes.

\begin{lstlisting}[caption=Input Field Accessibility, label=lst:input-focus]
// frontend/src/TooljetDatabase/Forms/ColumnForm.jsx
<input
  type="text"
  value={value}
  placeholder="Enter value"
  aria-label="Input field label"
  onChange={(e) => setValue(e.target.value)}
  onFocus={() => speak(value ? `Input field, current value: ${value}` : 'Input field')}
  disabled={isDisabled}
/>
\end{lstlisting}

Input fields utilize native HTML \texttt{<input>} elements which are inherently focusable without requiring explicit \texttt{tabIndex} attributes. Unlike interactive components that use \texttt{tabIndex=\{-1\}} to remove disabled elements from keyboard navigation, input fields rely on the native \texttt{disabled} attribute which automatically removes them from the tab sequence while also preventing user interaction. The \texttt{<input>} element implicitly provides \texttt{role="textbox"} to assistive technologies. Accessibility is ensured through \texttt{placeholder} for visual guidance, \texttt{aria-label} for screen reader identification, and \texttt{onFocus} callbacks that trigger screen reader announcements using the Web Speech API.

\textbf{Code Editor Component}

The code editor component utilizes the @uiw/react-codemirror library, which provides built-in keyboard accessibility for code editing within the canvas. For query panel code editors, we implemented Enter/Escape mode switching to enable keyboard navigation.

\begin{lstlisting}[caption=Code Editor Component, label=lst:codeeditor-focus]
// Canvas CodeMirror component
<CodeMirror
  value={value}
  placeholder={placeholder}
  theme={darkMode ? okaidia : githubLight}
  extensions={[languageExtension]}
  onChange={(value) => handleChange(value)}
  basicSetup={{
    lineNumbers: true,
    syntaxHighlighting: true,
    bracketMatching: true,
    autocompletion: true
  }}
  indentWithTab={true}
/>

// Query panel code editor with mode switching
const [isEditMode, setIsEditMode] = useState(false);

const handleWrapperKeyDown = (e) => {
  if (e.key === 'Enter' && !isEditMode) {
    e.preventDefault();
    setIsEditMode(true);
    speak('Entering edit mode');
    editorView.focus();
  }
};

const keymap = [
  {
    key: 'Escape',
    run: () => {
      if (isEditMode) {
        setIsEditMode(false);
        speak('Exiting edit mode');
        return true;
      }
    }
  },
  {
    key: 'Tab',
    run: () => {
      if (!isEditMode) return false; // Browser handles Tab navigation
      // Handle Tab indentation when in edit mode
    }
  }
];

<div
  tabIndex={isEditMode ? -1 : 0}
  onKeyDown={handleWrapperKeyDown}
  aria-label="Code editor. Press Enter to edit, Escape to exit."
/>
\end{lstlisting}

The CodeMirror library handles code editing shortcuts (Ctrl+Z, Ctrl+F) internally. Query panel editors add Enter to activate edit mode and Escape to exit, preventing Tab from indenting when navigating.

\textbf{Dynamic Component Properties}

Dynamic components within the app builder's properties panel are made focusable through systematic DOM traversal and attribute injection.

\begin{lstlisting}[caption=Dynamic Component Focusability, label=lst:dynamic-focus]
// frontend/src/_components/KeyboardNavigation/KeyboardNavigation.jsx
const makeElementsFocusable = () => {
  const items = document.querySelectorAll('.list-group-item, [role="button"]');
  
  items.forEach(item => {
    if (isElementVisible(item) && !item.hasAttribute('tabindex')) {
      item.setAttribute('tabindex', '0');
      item.addEventListener('keydown', (e) => {
        if (e.key === 'Enter' || e.key === ' ') {
          e.preventDefault();
          item.click();
        }
      });
    }
  });
};
\end{lstlisting}

This approach ensures that dynamically rendered components receive keyboard focus capabilities without requiring individual component modifications, maintaining consistency across the application.

\subsubsection{Keyboard-Based Interactions}

\textbf{Drag and Drop Keyboard Support}

Keyboard-based component placement provides an alternative to mouse-based drag-and-drop operations, enabling users to add components to the canvas using arrow keys for positioning and Enter for confirmation.

\begin{lstlisting}[caption=Keyboard Component Placement Implementation, label=lst:keyboard-dragdrop]
// frontend/src/AppBuilder/AppCanvas/KeyboardPlacementOverlay.jsx
const handleKeyDown = (e) => {
  if (!active) return;

  // Arrow keys - move the component preview
  if (['ArrowUp', 'ArrowDown', 'ArrowLeft', 'ArrowRight'].includes(e.key)) {
    e.preventDefault();
    moveKeyboardPlacementPosition(e.key);
  }
  // Enter - place component at current position
  else if (e.key === 'Enter') {
    e.preventDefault();
    handleDrop(component, 'canvas');
    cancelKeyboardPlacement();
  }
  // Escape - cancel placement
  else if (e.key === 'Escape') {
    e.preventDefault();
    cancelKeyboardPlacement();
  }
};
\end{lstlisting}

The keyboard placement system renders a visual overlay showing the component's preview position, which updates in real-time as users press arrow keys. This provides immediate visual feedback for precise positioning.

\textbf{Component Resizing Shortcuts}

Components support keyboard-based resizing through modifier keys combined with arrow keys, enabling precise dimensional adjustments without mouse interaction.

\begin{lstlisting}[caption=Keyboard Resizing Implementation, label=lst:keyboard-resize]
// frontend/src/AppBuilder/AppCanvas/WidgetWrapper.jsx
const handleKeyDown = (e) => {
  // = key for expanding, - key for shrinking
  if (e.key === '=') setResizeModifier('expand');
  else if (e.key === '-') setResizeModifier('shrink');
  
  // Arrow keys with active modifier
  if (isArrowKey && resizeModifier) {
    e.preventDefault();
    resizeComponentWithKeyboard(e.key, resizeModifier === 'expand');
    speak(`${resizeModifier === 'expand' ? 'Expanding' : 'Shrinking'} ${direction}`);
  }
};
\end{lstlisting}

The resizing mechanism prevents movement operations when a resize modifier is active, ensuring that arrow key presses modify dimensions rather than position.

\textbf{Component Deletion}

The Backspace key enables quick component removal with confirmation safeguards for destructive actions.

For Component Deletion, Undo/Redo, and Copy/Paste operations, we utilized the useHotkeys library with standard platform conventions:

\begin{lstlisting}[caption=Editor Keyboard Shortcuts, label=lst:keyboard-editor-shortcuts]
// frontend/src/Editor/Container.jsx
useHotkeys('backspace', () => removeMultipleComponents());
useHotkeys('control+z', () => handleUndo());
useHotkeys('control+shift+z', () => handleRedo());
useHotkeys('control+c', () => copyComponent());
useHotkeys('control+v', () => pasteComponent());

const handleUndo = () => {
  if (canUndo()) {
    undo();
    toast.info('Undone last action');
  }
};
\end{lstlisting}

These shortcuts support both single and multi-component operations while maintaining undo/redo compatibility. The undo/redo system maintains a history stack of application state changes, enabling users to revert or reapply modifications to components, properties, and layout configurations.

\textbf{Canvas Scroll Navigation}

Ctrl+Arrow key combinations control canvas viewport scrolling independent of component movement or focus changes.

\begin{lstlisting}[caption=Canvas Scroll Keyboard Control, label=lst:keyboard-scroll]
// frontend/src/_components/KeyboardNavigation/KeyboardNavigation.jsx
useHotkeys('control+up, control+down, control+left, control+right', (e) => {
  e.preventDefault();
  const canvas = document.querySelector('.real-canvas');
  const scrollAmount = 50;
  
  if (e.key === 'ArrowUp') canvas.scrollTop -= scrollAmount;
  else if (e.key === 'ArrowDown') canvas.scrollTop += scrollAmount;
  else if (e.key === 'ArrowLeft') canvas.scrollLeft -= scrollAmount;
  else if (e.key === 'ArrowRight') canvas.scrollLeft += scrollAmount;
  
  speak(`Scrolled canvas ${e.key.replace('Arrow', '')}`);
}, { enableOnTags: ['INPUT', 'TEXTAREA', 'SELECT'] });
\end{lstlisting}

The scroll navigation remains functional even when input fields have focus, utilizing the \texttt{enableOnTags} option to override default behavior for these specific shortcuts.

\subsubsection{Navigation Patterns}

\textbf{Arrow Key Navigation in Tab Options}

Tab-based interfaces within settings menus implement arrow key navigation for rapid movement between options, following ARIA authoring practices for tab panels.

\begin{lstlisting}[caption=Tab Options Arrow Key Navigation, label=lst:arrow-tabs]
// frontend/src/TooljetDatabase/Forms/RowForm.jsx
const handleTabKeyDown = (e) => {
  if (e.key === 'ArrowLeft' || e.key === 'ArrowRight') {
    e.preventDefault();
    const allTabs = Array.from(tablist.querySelectorAll('[role="button"]'));
    const currentIdx = allTabs.indexOf(e.currentTarget);
    const nextIdx = e.key === 'ArrowRight'
      ? (currentIdx + 1) % allTabs.length
      : (currentIdx - 1 + allTabs.length) % allTabs.length;
    allTabs[nextIdx]?.focus();
    speak(`${allTabs[nextIdx].textContent} tab`);
  }
};

<div role="tablist">
  <div role="button" tabIndex={activeTab === index ? 0 : -1} 
       aria-selected={activeTab === index} onKeyDown={handleTabKeyDown}/>
</div>
\end{lstlisting}

The tab navigation implements roving tabindex pattern, where only the active tab is in the tab sequence (\texttt{tabIndex=\{0\}}), while inactive tabs are removed (\texttt{tabIndex=\{-1\}}). Arrow keys move focus and update the roving tabindex.

\textbf{Dynamic Component Navigation}

The global keyboard navigation system automatically detects and integrates dynamically rendered components into the navigation sequence through continuous DOM observation.

\begin{lstlisting}[caption=Dynamic Component Navigation Integration, label=lst:dynamic-navigation]
// frontend/src/_components/KeyboardNavigation/KeyboardNavigation.jsx
const getFocusableElements = () => {
  const elements = document.querySelectorAll('.app-card, .sidebar-icon, [role="button"]');
  elements.forEach(el => {
    if (isElementVisible(el) && !el.hasAttribute('tabindex')) {
      el.setAttribute('tabindex', '0');
    }
  });
  return Array.from(elements);
};

// Monitor DOM changes
useEffect(() => {
  const observer = new MutationObserver(() => getFocusableElements());
  observer.observe(document.body, { childList: true, subtree: true });
  return () => observer.disconnect();
}, []);
\end{lstlisting}

The MutationObserver API monitors DOM changes, triggering updates to the focusable element list when components are added or removed, ensuring navigation sequences remain current.

\subsubsection{Focus Management}

\textbf{Focus Traps with ESC Key}

Modal dialogs and dropdown menus implement focus trapping to constrain keyboard navigation within the active context, preventing focus from escaping to background elements.

\begin{lstlisting}[caption=Menu Focus Trap Implementation, label=lst:focus-trap]
// frontend/src/TooljetDatabase/TableListItem/ActionsPopover/index.jsx
useEffect(() => {
  if (!show) return;
  
  const handleKeyDown = (e) => {
    if (e.key === 'Tab') {
      e.preventDefault();
      const items = getFocusableElements();
      const currentIndex = items.indexOf(document.activeElement);
      const nextIndex = e.shiftKey 
        ? (currentIndex - 1 + items.length) % items.length
        : (currentIndex + 1) % items.length;
      items[nextIndex]?.focus();
    } else if (e.key === 'Escape') {
      e.preventDefault();
      setShow(false);
      triggerButtonRef.current?.focus();
    }
  };
  
  document.addEventListener('keydown', handleKeyDown, true);
  return () => document.removeEventListener('keydown', handleKeyDown, true);
}, [show]);
\end{lstlisting}

Focus traps prevent Tab from moving focus outside the menu, cycling through menu items instead. Escape closes the menu and returns focus to the triggering element, maintaining user context.

\textbf{Visual Focus Indicators}

To provide clear visual feedback for keyboard navigation, we implemented CSS-based focus indicators using the \texttt{outline} property with a \textbf{2px thickness}. When an element receives focus via keyboard navigation, a visible outline appears around it, meeting WCAG 2.1 Level AA requirements for focus visibility.

\begin{lstlisting}[caption=Keyboard Focus Visual Indicators, label=lst:focus-indicator]
/* Focus outline - blue (2px thickness) */
button:focus:not([disabled]),
input:focus:not([disabled]),
[tabindex]:focus:not([tabindex="-1"]):not([disabled]),
.app-card:focus {
  outline: 2px solid #3E63DD;
  outline-offset: 2px;
  border-radius: 4px;
  box-shadow: 0 0 0 4px rgba(62, 99, 221, 0.15);
}

/* Dark mode - lighter blue outline (2px thickness) */
.dark-mode button:focus:not([disabled]),
.dark-mode [tabindex]:focus:not([tabindex="-1"]) {
  outline: 2px solid #5B9BFF;
  box-shadow: 0 0 0 4px rgba(91, 155, 255, 0.2);
}
\end{lstlisting}

The focus indicators use a consistent \textbf{2px solid outline} with \textbf{2px offset} from the element border, displaying a blue outline (\texttt{\#3E63DD}) for all focused elements. Dark mode employs a lighter blue (\texttt{\#5B9BFF}) to maintain sufficient contrast against darker backgrounds. A \textbf{4px box-shadow} with 15-20\% opacity provides additional visual emphasis, enhancing visibility for users with low vision.

\subsubsection{Keyboard Shortcuts Reference}

The following table summarizes all keyboard shortcuts implemented across the application's editor and app builder interfaces.

\begin{table}[h]
\centering
\caption{Implemented Keyboard Shortcuts}
\label{tab:keyboard-shortcuts}
\begin{tabular}{|l|l|p{6cm}|}
\hline
\textbf{Shortcut} & \textbf{Context} & \textbf{Function} \\
\hline
Arrow Keys & Component & Move selected component on canvas grid \\
\hline
= + Arrows & Component & Expand component dimensions \\
\hline
- + Arrows & Component & Shrink component dimensions \\
\hline
Backspace & Component & Delete selected component(s) \\
\hline
Ctrl + Z & Editor & Undo last action \\
\hline
Ctrl + Y & Editor & Redo last undone action \\
\hline
Ctrl + Shift + Z & Editor & Redo last undone action (alternative) \\
\hline
Ctrl + C & Component & Copy selected component \\
\hline
Ctrl + V & Canvas & Paste component from clipboard \\
\hline
Ctrl + X & Component & Cut selected component \\
\hline
Ctrl + D & Component & Duplicate selected component \\
\hline
Ctrl + Arrow & Canvas & Scroll canvas viewport \\
\hline
Tab & Universal & Navigate to next focusable element \\
\hline
Shift + Tab & Universal & Navigate to previous focusable element \\
\hline
Enter & Universal & Activate focused element / Toggle input mode \\
\hline
Space & Button/Toggle & Activate focused button or toggle \\
\hline
Escape & Menu/Modal & Close menu/modal, return to previous context \\
\hline
Arrow Left/Right & Tabs & Navigate between tab options \\
\hline
\end{tabular}
\end{table}

These shortcuts follow established accessibility patterns and platform conventions, ensuring intuitive operation for users familiar with keyboard-driven interfaces. All shortcuts remain functional when input fields have focus where contextually appropriate, utilizing event capture phase and selective \texttt{preventDefault()} calls to maintain expected behavior.
